\section{Estándares y Tecnologías Habilitadoras}

\subsection{3GPP Releases y Roadmap 6G}

Resulta revelador observar cómo 3GPP ha acelerado en los últimos años la incorporación de capacidades NTN dentro del ecosistema celular. 

La Release 17 estableció las bases con soporte básico para NTN en modo bent-pipe, mientras que la Release 18 introdujo mejoras sustanciales en corrección de frecuencia y timing advance. La Release 19 marca un punto de inflexión al incorporar payloads regenerativos y funcionalidades avanzadas de D2D. Análisis industriales recientes enfatizan que este camino prepara el terreno para una integración nativa TN-NTN en las primeras releases 6G \cite{Ericsson2025_NTN}. 

La documentación oficial de 3GPP confirma que Rel-19 ya contempla escenarios de movilidad mejorada y optimizaciones específicas para órbitas LEO \cite{3GPP_NTNR19}. No obstante, cabe matizar que el roadmap todavía presenta brechas importantes en lo que respecta a orquestación inteligente y soporte completo de D2D multi-banda.

\subsection{Técnicas de Mitigación Doppler y Movilidad}

Uno de los desafíos persistentes al operar en órbitas bajas radica en manejar velocidades relativas que generan desplazamientos Doppler de decenas de kHz. 

Las soluciones actuales combinan pre-compensación en el lado del satélite, estimación avanzada de canal en el dispositivo y mecanismos de timing advance mejorados. La Release 19 introduce procedimientos de handover TN-NTN más robustos que consideran predicción de trayectoria satelital. Estas técnicas resultan esenciales para mantener la sincronización y minimizar interrupciones en escenarios de movilidad continua.

\begin{figure}[t]
\centering
\includegraphics[width=1\linewidth]{fig4.pdf}
\caption{Protocolos de handover TN-NTN propuestos para convergencia D2D en redes 6G.}
\label{fig:handover_protocols}
\end{figure}

Aunque estas mejoras representan un avance significativo, en muchos casos estudiados todavía generan overhead de señalización que impacta el consumo energético de dispositivos de bajo perfil.

\subsection{Gestión de Recursos e Interferencias}

Lejos de ser un asunto resuelto, la gestión eficiente de recursos en entornos de convergencia TN-NTN exige enfoques más sofisticados que los utilizados en redes puramente terrestres. 

La coexistencia de enlaces terrestres y satelitales genera patrones complejos de interferencia, especialmente en bandas compartidas. Encuestas recientes sobre D2D satelital destacan la importancia de técnicas como beamforming dinámico, control de potencia adaptativo y asignación de recursos asistida por IA para mitigar estas interferencias \cite{Pasandi2024_D2DSurvey}.

\begin{table}[t]
\centering
\caption{Requisitos técnicos por banda para D2D-NTN}
\label{tab:band_requirements}
\begin{tabular}{lccc}
\hline
Banda & Ancho de Banda & Doppler Máx. (kHz) & Consumo Energético Relativo \\
\hline
Sub-7 GHz (n256) & 10-20 MHz & $\sim 15$ & Bajo \\
Ka-band & 100-400 MHz & $\sim 45$ & Alto \\
\hline
\end{tabular}
\end{table}

La Tabla 3 resume los principales requisitos y trade-offs por banda. La arquitectura propuesta incorpora un módulo de orquestación que realiza asignación conjunta de recursos, priorizando enlaces terrestres en áreas densas y satelitales en zonas de baja cobertura. Esta aproximación mejora notablemente la eficiencia espectral global, aunque aumenta la complejidad de implementación en el core network.